\begin{lemma}[Two-color components have a spanning path]
Let $\varphi$ be proper.  For every node $z$ of the two-color incidence
graph, there are nodes $u,v$ and a simple path from $u$ to $v$ whose set of
nodes is exactly the $(\alpha,\beta)$-component through $z$.
\end{lemma}
\begin{proof}
Let $K$ be the component from \noderef{TwoColorComponent}.  It is finite
and nonempty, since it contains $z$.  Among all simple paths in the
incidence graph whose nodes lie in $K$, choose one, say $P$, with the
largest possible number of nodes.  Such a path exists: the length-zero path
at $z$ is admissible, and the finite set $K$ bounds all possible support
sizes.

Suppose that a node $x\in K$ is not on $P$.  Since $K$ is a connected
component, there is a walk from $x$ to a node of $P$.  Delete repetitions
from a shortest such walk, and stop at its first node $y$ on $P$.  Its
predecessor $x'$ is outside $P$ and is adjacent to $y$.

If $y$ is an end of $P$, adjoining the edge $x'y$ and the preceding simple
part of the chosen walk extends $P$ to a simple path with more nodes, which
contradicts maximality.  If $y$ is internal to $P$, it already has two
distinct neighbors on $P$; together with the distinct outside neighbor
$x'$, it has degree at least three.  This contradicts
\noderef{TwoColorIncidenceDegreeAtMostTwo}.  Therefore no such $x$ exists.
Every node of $K$ lies on $P$, while $P$ was chosen inside $K$, so the
support of $P$ is exactly $K$.
\end{proof}
